% Generated draft: requires mathematical and source review.
\documentclass[11pt]{article}
\usepackage[margin=25mm]{geometry}
\usepackage{fontspec}
\setmainfont{Latin Modern Roman}
\usepackage{amsmath,amssymb,mathtools,unicode-math}
\setmathfont{Latin Modern Math}
\usepackage{enumitem,needspace,xurl,hyperref}
\setlength{\parindent}{0pt}
\setlength{\parskip}{5pt}
\setlength{\emergencystretch}{3em}
\newcommand{\N}{\mathbb{N}}
\newcommand{\Z}{\mathbb{Z}}
\newcommand{\Q}{\mathbb{Q}}
\newcommand{\R}{\mathbb{R}}
\newcommand{\C}{\mathbb{C}}
\newcommand{\F}{\mathbb{F}}
\newcommand{\A}{\mathbb{A}}
\newcommand{\PP}{\mathbb{P}}
\newcommand{\E}{\mathbb{E}}
\newcommand{\eps}{\varepsilon}
\newcommand{\dd}{\,\mathrm{d}}
\begin{document}
\section*{UnsolvedMath ID: 800300}
\textbf{Source problem code:} catalogue entry
\subsection{Definitions and mathematical statement}
Let \(S^1=\R/\Z\) be the circle, equipped with normalized Lebesgue measure \(m\), so that \(m(S^1)=1\). Let \(G\) be a finitely generated group of \(C^2\) diffeomorphisms of \(S^1\). Its action is \emph{minimal} if every orbit is dense:
\[
\overline{\{g(x):g\in G\}}=S^1\qquad\text{for every }x\in S^1.
\]
Every such diffeomorphism preserves the Lebesgue measure class, although it need not preserve Lebesgue measure itself. Here, ergodicity with respect to Lebesgue measure is interpreted in the nonsingular sense: every Lebesgue measurable set \(A\subseteq S^1\) satisfying
\[
m\bigl(g^{-1}(A)\mathbin{\triangle}A\bigr)=0\qquad\text{for every }g\in G
\]
has \(m(A)\in\{0,1\}\), where \(\triangle\) denotes symmetric difference.

\textbf{Question (Navas, 2017).} Must every minimal action of a finitely generated group of \(C^2\) circle diffeomorphisms be ergodic with respect to Lebesgue measure in this sense?

\textbf{Unresolved source qualifications.} The supplied statement does not specify whether all group elements must preserve orientation. This convention requires confirmation. It also does not explicitly define ergodicity; the formulation above uses nonsingular ergodicity and introduces no assumption that the group preserves Lebesgue measure.
\subsection{Short English statement}
Navas asks whether minimality of a finitely generated group action by \(C^2\) circle diffeomorphisms forces every invariant measurable set to have zero or full Lebesgue measure. The orientation convention and the intended definition of ergodicity require source confirmation.
\subsection{Sources}
\begin{itemize}
\item Local UnsolvedMath registry, supplied complete record: permanent ID \(800300\), source problem code \texttt{catalogue entry}, title \emph{Group Actions on 1-Manifolds --- Question 17}. The registry attributes the question to Navas (2017) and labels it open; this is supplied status metadata, not an independently verified assessment.
\item Navas, \emph{Group Actions on 1-Manifolds: A List of Very Concrete Open Questions} (2017), \url{https://arxiv.org/abs/1712.06462}. The supplied provenance identifies the source item as \texttt{Question 17 (\#16)}, records access on 2026-10-16, and gives the extraction method as \texttt{arxiv\_latex}. These provenance details are supplied metadata; the source was not independently consulted for this draft.
\end{itemize}
\end{document}
